% Author: Keshav
\section{Introduction}
\label{sec:intro}

Entity resolution decides which records in different databases describe the same real-world
entity. Our task comes from the Amazon ML Challenge 2026: Source~1
(S1) is a reference list of businesses, and every record of Sources~2 and~3 (S2, S3) has to be linked to the
S1 entity it describes, or to none. The data covers the US and India. The noise is varied:
\begin{itemize}
  \item abbreviated legal forms (\emph{private limited} / \emph{pvt ltd});
  \item phonetic romanisations of Indian words (\emph{praivtt}, \emph{enjiniyring});
  \item names written in Devanagari or Malayalam, while S1 is always in Latin script;
  \item missing addresses;
  \item near-duplicate branches of one chain that differ only in a house number.
\end{itemize}
About a quarter of S2/S3 records have no S1 entity at all. The metric is macro \fhalf{} per S1 entity,
which weighs precision more than recall, so a wrong link costs more than a missed one.

This makes the task a good test bed for the course material. It needs tokenization that copes with
three scripts, string similarity and edit distance, sequence labelling (address parsing), recurrent
networks, pre-trained transformers with full and parameter-efficient fine-tuning, and Mixture-of-Experts
models. We built one shared pipeline (\S\ref{sec:overview}) in which every matcher scores \emph{the same}
candidate pairs with \emph{the same} decision rule and threshold protocol, so that their results are
directly comparable. The sections below follow the pipeline: data and splits (\S\ref{sec:data}), tokenization and retrieval (\S\ref{sec:tok}), parsing, the classical, recurrent and transformer matchers including Ditto (\S\ref{sec:parsing}--\ref{sec:ce}), Mixture-of-Experts fusion (\S\ref{sec:fusion}) and a per-script error analysis (\S\ref{sec:script}).
Table~\ref{tab:main} in \S\ref{sec:fusion} collects all results.
